\documentclass[10pt]{article}
% Math Packages
\usepackage{amsmath, mathtools}
\usepackage{amssymb}
\usepackage{amsthm}
\usepackage{amsfonts}
\usepackage{bbm}
\usepackage{breqn}
\usepackage[margin=1in]{geometry}
\usepackage{graphicx}
\usepackage{tikz}
\usetikzlibrary{arrows.meta}
\usetikzlibrary{calc}
\usepackage{forest}
\usepackage{tikz-qtree}
\graphicspath{ {./images/} }
\usepackage{hyperref}
\usepackage[capitalize]{cleveref}
\usepackage[shortlabels]{enumitem}
\usetikzlibrary{arrows,matrix,positioning}
\usepackage{multicol}

% for the pipe symbol
\usepackage[T1]{fontenc}

% Citing theorems by name. (source: https://tex.stackexchange.com/questions/109843/cleveref-and-named-theorems)
\makeatletter
\newcommand{\ncref}[1]{\cref{#1}\mynameref{#1}{\csname r@#1\endcsname}}

\def\mynameref#1#2{%
  \begingroup
  \edef\@mytxt{#2}%
  \edef\@mytst{\expandafter\@thirdoffive\@mytxt}%
  \ifx\@mytst\empty\else
  \space(\nameref{#1})\fi
  \endgroup
}
\makeatother

% Colorful Notes
\usepackage{color} \definecolor{Red}{rgb}{1,0,0} \definecolor{Blue}{rgb}{0,0,1}
\definecolor{Purple}{rgb}{.5,0,.5} \def\red{\color{Red}} \def\blue{\color{Blue}}
\def\gray{\color{gray}} \def\purple{\color{Purple}}
\newcommand{\rnote}[1]{{\red [#1]}} % \rnote{foo} gives '[foo]' in red
\newcommand{\pnote}[1]{{\purple [#1]}} % \pnote{foo} gives '[foo]' in purple
\newcommand{\bnote}[1]{{\blue #1}} % \bnote{foo} gives 'foo' in blue
\newcommand{\gnote}[1]{{\gray #1}} % \gnote{foo} gives 'foo' in gray
\newcommand{\Max}[1]{{\purple [#1]}} % \bnote{foo} then 'foo' is blue


% Claim numbering (the counter restarts after each proof environment)
\newcounter{claimcount}
\setcounter{claimcount}{0}
\newenvironment{claim}{\refstepcounter{claimcount}\par\addvspace{\medskipamount}\noindent\textbf{Claim \arabic{claimcount}:}}{}
\usepackage{etoolbox}
\AtBeginEnvironment{proof}{\setcounter{claimcount}{0}}
\newenvironment{claimproof}{\par\addvspace{\medskipamount}\noindent\textit{Proof of Claim  \arabic{claimcount}.}}{\hfill\ensuremath{\qedsymbol} \tiny{Claim}

  \medskip}
% Add claim support to cleverref
\crefname{claimcount}{Claim}{Claims}


% Math Environments
\newtheorem{theorem}{Theorem}
\newtheorem{assumption}[theorem]{Assumption}
\newtheorem{lemma}[theorem]{Lemma}
\newtheorem{proposition}[theorem]{Proposition}
\newtheorem{corollary}[theorem]{Corollary}
\newtheorem{question}[theorem]{Question}
\theoremstyle{definition}
\newtheorem{definition}[theorem]{Definition}
\newtheorem{remark}[theorem]{Remark}
\newtheorem{example}[theorem]{Example}
\newtheorem{notation}[theorem]{Notation}
\newtheorem{problem}[theorem]{Problem}

% Matrices and Column Vectors. 
\usepackage{stackengine}
\setstackgap{L}{1.0\normalbaselineskip}
\usepackage{tabstackengine}
\setstackEOL{;}% row separator
\setstackTAB{,}% column separator
\setstacktabbedgap{1ex}% inter-column gap
\setstackgap{L}{1.5\normalbaselineskip}% inter-row baselineskip
\let\nmatrix\bracketMatrixstack  %Usage: \nmatrix{1,2,3\4,5,6}
\newcommand\cv[1]{\setstackEOL{,}\bracketMatrixstack{#1}} %usage: \cv{1,2,3}

% Custom Math Coqmmands
\newcommand{\vt}{\vskip 5mm} % vertical space
\newcommand{\fl}{\noindent\textbf} % first line
\newcommand{\Fl}{\vt\noindent\textbf} % first line with space above
\newcommand{\norm}[1]{\left\lVert#1\right\rVert} % norm
\newcommand{\pnorm}[1]{\left\lVert#1\right\rVert_p} % p-norm
\newcommand{\qnorm}[1]{\left\lVert#1\right\rVert_q} % q-norm
\newcommand{\1}[1]{\textbf{1}_{\left[#1\right]}} % indicator function 
\def\limn{\lim_{n\to\infty}} % shortcut for lim as n-> infinity
\def\sumn{\sum_{n=1}^{\infty}} % shortcut for sum from n=1 to infinity
\def\sumkn{\sum_{k=1}^{n}} % shortcut for sum from k=1 to n
\def\sumin{\sum_{i=1}^{n}} % shortcut for sum from i=1 to n
\def\SAs{\sigma\text{-algebras}} % shortcut for $\sigma$-algebras
\def\SA{\sigma\text{-algebra}} % shortcut for $\sigma$-algebra
\def\Ft{\mathcal{F}_t} % time-indexed sigma-algebra (t)
\def\Fs{\mathcal{F}_s} % time-indexed sigma-algebra (s)
\def\F{\mathcal{F}} % sigma-algebra
\def\G{\mathcal{G}} % sigma-algebra
\def\R{\mathbb{R}} % Real numbers
\def\N{\mathbb{N}} % Natural numbers
\def\Z{\mathbb{Z}} % Integers
\def\E{\mathbb{E}} % Expectation
\def\P{\mathbb{P}} % Probability
\def\Q{\mathbb{Q}} % Q probability
\def\dist{\text{dist}} %Text 'dist' for things like 'dist(x,y)'
\newcommand{\indep}{\perp \!\!\! \perp}  %independence symbol
\def\Var{\mathrm{Var}} % Variance
\def\tr{\mathrm{tr}} % trace

% Brackets and Parentheses
\def\[{\left [}
    \def\]{\right ]}
% \def\({\left (}
%   \def\){\right )}



\usepackage{color}
\definecolor{Red}{rgb}{1,0,0}
\definecolor{Blue}{rgb}{0,0,1}
\definecolor{Purple}{rgb}{.5,0,.5}
\def\red{\color{Red}}
\def\blue{\color{Blue}}
\def\gray{\color{gray}}
\def\purple{\color{Purple}}
\definecolor{RoyalBlue}{cmyk}{1, 0.50, 0, 0}
\newcommand{\dempfcolor}[1]{{\color{RoyalBlue}#1}} 
\newcommand{\demph}[1]{\dempfcolor{{\sl #1}}}

% comment exactly one of the following line to show / hide the solutions
% \newcommand{\solution}[1]{{\purple #1}} % uncomment to show the solutions
\newcommand{\solution}[1]{} % uncomment to hide the solutions



\title{Lecture Notes for Math 372: \\Elementary Probability and Statistics}
\date{Last updated: \today}
% \author{mh}

\begin{document}
\begin{center}
  \section*{Math 372: Homework 03}
  \textit{Due Wednesday, Feb 4}
\end{center}

\begin{problem}
  The year is 1532, and 31-year old Girolamo Cardano has been appointed mathematics chair at the
  University of Milan. In celebration, Cardano decides to host a dinner party. In his cellar he finds
  30 bottles of wine: 8 bottles of zinfandel, 10 bottles of merlot, and 12 bottles of cabernet, all
  from different wineries. (Mr.~Cardano only drinks red wine).
  \begin{enumerate}[(a),itemsep=0em]
    \item He wants to serve 3 bottles of zinfandel and the serving order is important, how many ways are
    there to do this?
    \item Suppose 6 bottles are randomly selected. How many ways are there to do this?
    \item If 6 bottles are randomly selected, how many ways are there to obtain two bottles of each
    variety?
    \item What's the probability of that happening?
    \item If 6 bottles are randomly selected, what's the probability that they are all of the same
    variety (i.e., all zinfandel, all merlot, or all cabernet)?
  \end{enumerate}
\end{problem}




\begin{problem}[Exercise 2.3.34]
  Computer keyboard failures can be attributed to electrical defects or
  mechanical defects. A repair facility currently has 25 failed keyboards, 6
  of which have electrical defects and 19 of which have mechanical defects.
  \begin{enumerate}[(a),itemsep=0em]
    \item How many ways are there to randomly select 5 of these keyboards for
    a thorough inspection (without regard to order)?
    \item In how many ways can a sample of 5 keyboards be selected so that
    exactly 2 have an electrical defect?
    \item If a sample of 5 keyboards is randomly selected, what is the
    probability that at least 4 of these will have a mechanical defect?
  \end{enumerate}
\end{problem}




\begin{problem}
  We return to the dice-rolling pirate, Diego. Suppose Diego rolls two dice and then adds the numbers
  together. The possible outcomes are $2,3,4,5,6,7,8,9,10,11, \text{ and }12$, but recall that these
  numbers are not all equally likely. 
  \begin{enumerate}[(a),itemsep=0em]
    \item What is the probability that the sum of the two dice is greater than 4? 
    \item Given that the sum of the dice is greater than 4, what is the
    probability that the dice added up to 7?
    \item Diego rolls his dice again. What is the probability that the dice add up to an even number?
    \item Diego rolls his dice once more. What is the probability that the dice add up to an odd
    number?
    \item Diego rolls his dice another time. Given that the dice add up to an even number, what is the
    probability that Diego rolled at least one 4?
    \item Diego rolls his dice again, and the two dice sum to 7. What is the probability that he rolled
    a 3 or a 5?
    \item Diego rolls his dice again. What is the probability that at least one of Diego's dice was a
    3?
    \item Given that at least one of Diego's dice was a 3, what is the probability that his dice add up
    to 8 or 9?
    \item Diego decides to roll his pair of dice another time. What is the
    probability that they sum to 4 or 7?
    \item As one last experiment Diego decides to keep rolling the two dice until he gets a sum of
    either 4 or a 7, at which point he will stop and go back to his ship. What's the probability that
    his last roll was 4?
  \end{enumerate}
\end{problem}

\begin{problem}
  A vampire goes to the blood bank looking to find some type O+ blood (the most delicious type). He finds
  8 unlabeled bags of blood. Only one of the bags is O+, but the vampire doesn't know which one, so he
  resorts to taste-testing to find the O+ bag.
  \begin{enumerate}[(a),itemsep=0em]
    \item  What is the probability that he must test at
    least 5 bags to find the desired type?
    % \noindent \textit{Solution:} Let $X$ be the number of tested bags. We want to find $\P\left[X \geq
    %   3\right]$.
    % Let $A= \left[ \text{first bag isn't O+} \right]$. Let $B = \left[ \text{second bag isn't O+} \right]$
    % \begin{align*}
    %   \P\left[ X \geq 3\right] 
    %   &= \P\left[A\cap B \right]\\
    %   &= \P\left[A \right] \P\left[B\mid A \right]\\
    %   &= \frac{3}{4}\cdot \frac{2}{3}\\
    %   &= \frac{1}{2}.
    % \end{align*}
    \item What's the probability that he must test exactly $5$ bags?
    % \noindent \textit{Solution:} We want $\P\left[X=3 \right]=\P\left[\text{third bag is O+} \right]$.
    % \begin{align*}
    %   \P\left[\text{third bag is O+} \right] 
    %   &= \P\left[\text{third bag is}\mid  \text{first isn't} \cap
    %     \text{second isn't}  \right] \times \P\left[\text{second isn't}\mid \text{first isn't} \right] \cdot
    %   \P\left[\text{first isn't} \right]\\
    %   &= \frac{1}{2}\cdot \frac{2}{3}\cdot \frac{3}{4}\\
    %   &= \frac{1}{4}
    % \end{align*}
  \end{enumerate}
\end{problem}


\begin{problem}
  A 5-card poker hand is a \textbf{full house} it has a 3-of-a-kind and 2-of-a-kind.
  \begin{enumerate}[(a),itemsep=0em]
    \item How many possible full houses are there?
    % The number of possible full houses
    % \begin{equation*}
    %   {13\choose 1}\times {4\choose 3}\times {12\choose 1}\times {4\choose 2} = 3,744
    % \end{equation*}
    \item Suppose you are dealt a random poker hand. What's the probability
    that you get a full house?
    % The probability of getting a full house is
    % \begin{equation*}
    %   \P\left[\text{full house} \right] = \frac{3,744}{2,598,960} = 0.00144
    % \end{equation*}
  \end{enumerate}
\end{problem}

\end{document}